\documentclass[a4paper]{article}

\usepackage[default]{lato}
\usepackage[swedish]{babel}
\usepackage[dvipsnames]{xcolor}
\usepackage[T1]{fontenc}
\usepackage[utf8]{inputenc}
\usepackage{graphicx}
\usepackage{fancyhdr}
\usepackage{lastpage}
\usepackage{parskip}
\usepackage[colorlinks=true,citecolor=black,linkcolor=black]{hyperref}
\usepackage[style=ieee]{biblatex}

\addbibresource{bib.bib}

\definecolor{cerise}{RGB}{232, 61, 132}
\definecolor{gray3}{RGB}{66, 66, 66}

\addtolength{\textheight}{4cm}
\addtolength{\voffset}{-2cm}
\addtolength{\textwidth}{2cm}
\addtolength{\hoffset}{-1cm}
\setlength{\headheight}{40pt}

\pagestyle{fancy}
\thispagestyle{fancy}
\fancyhead[L]{\small{
    Motion till Revisions-SM
\\\textcolor{gray3}{
    2023-03-28
}}}
\fancyhead[C]{\includegraphics[width=1.25cm,height=1.25cm]{./skold-color.pdf}}
\fancyhead[R]{\small{\textcolor{gray3}{
    Mathias Magnusson\\ % Skriv ditt namn här
    sida \thepage{} av \pageref*{LastPage}}}}
\fancyfoot[C]{\small{\textcolor{gray3}{Konglig Datasektionen, THS 100 44, datasektionen.se}}}
\renewcommand{\headrulewidth}{0pt}

\title{\textcolor{cerise}{\textbf{
    Motion angående\\
    dJulkalendern någonting någonting nämnd\\
    eller\\
    hur vi ser till att dJulkalendern 2023 inte blir av
}}}

\date{}

\begin{document}

\maketitle
\thispagestyle{fancy} % Härligt finhack för att få header på första sidan

\section*{\textcolor{cerise}{
    Bakgrund
}}

dJuvliga hälsningar Datasektionen!

dJulkalendern hölls för första gången 2014 \cite{lol} då Ludwig Janiuk var \textit{eldsjälen i
kärnan} och jobbet skedde \textit{måttligt planerat och måttligt strukturerat}
\cite{weigelt_20}.

Efter Budget-SM år 2020 blev dJulkalendern tack vare Christian ett faktiskt projekt under Konglig
Datasektionen och så har det nu varit i tre år. Det har fungerat bra och så skulle vi för alla del
kunna fortsätta ha det. Däremot vore det mer passande att dJulkalendern istället skulle vara en
nämnd. Bland annat skulle det glädja Dr. Ek, men framförallt skulle det underlätta för mer
kontinuerligt arbete på kalendern under årets gång.

Dessutom har dHöstkalendern anordnats under mottagningen av projektgruppen för dJulkalendern men,
den har skett (och förberetts) innan projektet ens haft någon verksamhetsplan eller budget
\cite{kaffe_dm_22}.

\vspace{1cm}

\begin{center}
    \includegraphics[width=0.5\textwidth]{./djulutan23.jpg}
\end{center}

\clearpage

\section*{\textcolor{cerise}{
    Förslag till beslut
}}

	Mot bakgrund av ovanstående yrkar jag:

\newcounter{attnummer}
\setcounter{attnummer}{1}
\begin{list}{\bf att$_{\theattnummer}$\stepcounter{attnummer}}{}
    \item i reglementet lägga till ``§3.22 dJulkalendern''

    \item i reglementet lägga till ``§3.22.1 Ändamål'' med följande innehåll:
        \begin{quote}
            dJulkalenderns ändamål är att främja intresse för etisk hackning, problem-lösning och
            CTF-tävlingar inom Konglig Datasektionen.
        \end{quote}

    \item i reglementet lägga till ``§3.22.2 Organisation'' med följande innehåll:
        \begin{quote}
            dJulkalendern leds av Tomtemor/-far.

            Övriga medlemmar utses av Tomtemor/-far.
        \end{quote}

    \item i reglementet lägga till ``§3.22.3 Verksamhet'' med följande innehåll:
        \begin{quote}
            dJulkalendern ska varje december anordna en CTF-liknande tävling på den världsvida
            webben för Konglig Datasektionsns medlemmar samt övriga intresserade.

            dJulkalendern ska även anordna dHöstkalendern, en mindre variant av dJulkalendern för
            nyantagna medlemmar till Konglig Datasektionen under mottagningen.
        \end{quote}

    \item i reglementet lägga till ``§5.2.22 Tomtemor/-far'' med följande innehåll:
        \begin{quote}
            Är ordförande för dJulkalendern.

            Väljs på Val-SM. Har läsår som mandatperiod.
        \end{quote}

    \item för dJulkalendern anta den bifogade \hyperref[fig:budget]{budgeten}.

    \item öppna upp för fyllnadsval till Tomtemor/-far.
\end{list}

\vspace{1cm}

Mathias Magnusson, Tomtefar \\
dAnkan, dAnkan

\printbibliography


\begin{figure}
    \includegraphics[width=1\textwidth]{./budget.pdf}
    \label{fig:budget}
\end{figure}

\end{document}
